% ============================================================
% 09conclusiones.tex
% Conclusiones finales
% ============================================================

\section{Conclusiones}
\label{sec:conclusiones}

Este trabajo estimó la probabilidad de victoria en el Pok\'emon Trading Card
Game a partir de un solo estado de la partida. Cada partida quedó completa
en un único subconjunto, y el conjunto de prueba se usó una sola vez. El
modelo principal, potenciación del gradiente con 62 características
públicas, obtuvo un AUC de 0.833 en partidas nuevas y cerca de 0.79 con
días posteriores y arquetipos de mazo nuevos. Superó a las cuatro
referencias en los tres escenarios, y en su propia población sus
probabilidades ya están calibradas.

La predicción mejora a medida que avanza la partida. Al inicio es modesta,
y en la población conocida no supera a un modelo que solo sabe el turno,
quién empezó y qué Pok\'emon están activos. Desde la etapa media domina el
estado del tablero, y al final el AUC supera 0.91 en todos los escenarios.
Con días o arquetipos nuevos, conocer los Pok\'emon activos deja de servir y
el estado sigue funcionando. Esa es la propiedad que necesita un evaluador
de posiciones.

Con el paso del tiempo, el modelo empeoró más en ordenar posiciones que en
calibrar sus probabilidades, al contrario de lo reportado en otros
dominios. Recalibrar con un día intermedio mejoró a B4, pero empeoró a B3,
así que no es una solución general. SHAP mostró que el modelo decide por
premios, descartes y mazo, sin señales de fuga, e identificó una
característica que falta: los premios que entrega cada Pok\'emon.

El modelo final, B4 entrenado en el escenario A y sin calibrador, se integró
en un agente con cuatro roles. Un modelo de lenguaje planifica y redacta,
pero toda cifra sale del modelo y un verificador determinista revisa cada
respuesta antes de mostrarla. En la validación, las 31 respuestas finales
tuvieron cifras trazables, el agente se negó a dar información oculta y no
siguió instrucciones inyectadas en los datos. Un simulador gráfico permite
a cualquier persona armar un estado y ver cómo cambia la probabilidad. El
límite principal es del modelo: en las etapas media y final, la mayoría de
sus errores muy seguros llegan sin advertencia, así que sus probabilidades
deben leerse como estimaciones, nunca como certezas.

Como trabajo futuro quedan repetir los experimentos con varias semillas,
agregar la característica de premios por Pok\'emon, entrenar la variante
de 64 características, seguir midiendo el desempeño a medida que cambie la
población de agentes y mostrar junto a cada probabilidad muy segura cuántas
veces fallan esas predicciones en su etapa.